% Declare the start of the appendices (this changes chapter numbering to letters)
\appendix

\chapter{ESP32 Data Acquisition}

This appendix contains the C++ firmware developed for the ESP32-WROOM-32D microcontroller. The software is responsible for the high-speed data acquisition stage of the stethoscope calibration rig.

\section{Timer Interrupt and Serial Transmission}

The following script initializes a hardware timer interrupt to trigger a 12-bit Analog-to-Digital (ADC) conversion exactly every 50 microseconds. This guarantees a strict 20,000 Hz sampling rate. The acquired integer values are subsequently pushed to the serial buffer at a baud rate of 500,000 bps to prevent data bottlenecking.

\begin{lstlisting}[language=C++, caption={ESP32 firmware for 20 kHz hardware timer ADC sampling and high-speed UART transmission.}, label={lst:esp32_firmware}]
#include <Arduino.h>

// Define hardware pins
const int adcPin = 34; // Connect to MAX9814 AC-coupled output
hw_timer_t * timer = NULL;
volatile bool sampleFlag = false;

// Interrupt Service Routine (ISR) triggered every 50us
void IRAM_ATTR onTimer() {
  sampleFlag = true;
}

void setup() {
  // Initialize UART at 500,000 baud
  Serial.begin(500000);
  
  // Set ADC resolution to 12-bit (0-4095)
  analogReadResolution(12);
  
  // Configure hardware timer (Prescaler 80 = 1us ticks)
  timer = timerBegin(0, 80, true);
  timerAttachInterrupt(timer, &onTimer, true);
  
  // Set alarm to trigger every 50 microseconds (20 kHz sampling)
  timerAlarmWrite(timer, 50, true);
  timerAlarmEnable(timer);
}

void loop() {
  if (sampleFlag) {
    sampleFlag = false;
    // Read ADC and transmit as raw integer
    int adcValue = analogRead(adcPin);
    Serial.println(adcValue);
  }
}
\end{lstlisting}